% Verbatim copy of /Users/alingalatan/workspaces/grothe/handoff/manuscript/legacy/lower-bound/sixth-moment-bridge.tex (file dated 2026-10-25). Only this line was added.
\section{The sixth-moment tensor bridge and the Gaussian vector witness}
\label{sec:sixth-tensor-bridge-audit}
\chapsummary{The tensor premises of the scalar constant $C_6=40.45313$
are proved for genuine complex Wick cubics. This is \st{P,C}, not a new
Lean theorem. The resulting radius is $0.969508314$. An exact vector
witness supplies the other, vector-valued transfer input on paper;
the dimension-free scalar cap remains the numerical obstruction.}

\paragraph{Provenance and formalization boundary.}
The 21 October note \emph{Audit and proofs for complex Wick cubics}
independently derives the identities below. The same mathematical chain
is archived in \pathref{handoff/lower-v4/proofs/} and
\pathref{handoff/lower-v5/local/}, especially
\code{MIXED\_CONTRACTION\_CAUCHY.md},
\code{CHARGE\_SPLIT\_SQUARE\_IMPROVEMENT.md}, and
\code{CHARGE\_SPLIT\_SCALAR\_CERTIFICATE.md}.
The audit replay checks the contraction identities, both sums of squares,
and the archived rational scalar certificate. No tensor premise remains
an unproved mathematical assumption for the class specified here.
Its construction and connection to \code{SixthMoment.cumulant\_bound}
are still absent from the supplied Lean interface.

\subsection{A scale-invariant moment theorem}
\begin{theorem}[\st{P,C}: genuine bidegree $(2,1)$]
\label{thm:sixth-tensor-bridge-audit}
For every finite-dimensional $q\in\mathcal W_{2,1}$, let
$e=\E|q|^2$ and $\kappa=\E|q|^4-2e^2$. Then
\[
 0\le\kappa\le42e^2,\qquad
 \E|q|^6\le9e\E|q|^4-12e^3+
             \frac{4045313}{100000}\kappa^{3/2}.
\]
The first assertion follows from Theorem~\ref{thm:sharp-influence-audit}.
The second does not hold for a general isotropic pair of real cubic
chaoses without the bidegree restriction.
\end{theorem}

\subsection{The two finite-sum identities}
Let $f$ be the real symmetric kernel of $\sqrt2\Re q=I_3(f)$,
$F_{ab,c}=f_{abc}$, $G=F^*F$, $C=FF^*$ and $S=\operatorname{Sym}_4 C$.
Use full ordered-index Hilbert tensor norms. Put
\[
 s=\|G\|^2,\quad h=\|S\|^2,\quad
 \eta=\|G_{\rm anti}\|^2,\quad \tau=\|S_{+2,+2}\|^2.
\]
Here $h$ is a tensor norm, not the directional influence.
In an orthonormal isotropic basis write $A_{ij,k}$ for the $(+,+,-)$
entries of $f$. They are symmetric in $i,j$; the opposite-charge
entries are their conjugates. Define
\[
\begin{aligned}
 P_{ij}&=\sum_{ab}A_{ia,b}\overline{A_{ja,b}},&
 Q_{ij}&=\sum_{ab}A_{ab,i}\overline{A_{ab,j}},&
 R_{ij}&=\sum_{ab}A_{ia,b}A_{jb,a},\\
 X_{ij;kl}&=\sum_a A_{ia,k}\overline{A_{la,j}},&
 Y_{ij;kl}&=\sum_a A_{ij,a}\overline{A_{kl,a}},&
 K_{ijk;l}&=\sum_a A_{ij,a}A_{ak,l}.
\end{aligned}
\]
Put $p_0=\|P\|^2$, $q_0=\|Q\|^2$, $v_0=\|R\|^2$,
$r_0=\langle P,\overline Q\rangle=\|K\|^2$ and
$T_0=\langle X,HX\rangle$, where $H$ swaps the first two indices of
$X$. Let $V$ swap its last two indices, $X_{\pm\pm}=(I\pm H)(I\pm V)X/4$,
and let $\pi$ cyclically permute the first three indices of $K$.
Relabeling finite sums gives
\[
\begin{gathered}
 \|X\|^2=p_0,\quad\|Y\|^2=q_0,\quad
 \langle X,VX\rangle=\langle X,Y\rangle
 =\langle K,\pi K\rangle=T_0,\quad
 \langle X,HVX\rangle=v_0,\\
 s=8p_0+2q_0+8r_0+8v_0,\qquad
 h=(s+48T_0)/3,\qquad \eta=8v_0,\\
 \tau=(4p_0+q_0+4v_0+16T_0)/9.
\end{gathered}
\]
Consequently the following are exact polynomial identities:
\begin{align}
37s-39h
 &=720\|X_{--}\|^2+48\|Y-X_{++}\|^2
                       +96\|K-\pi K\|^2,\label{eq:c6-sos-one}\\
216\tau-27h+5s-8\eta
 &=24\|X-HX\|^2+16\|P-\overline Q\|^2.
 \label{eq:c6-sos-two}
\end{align}
They prove $h\le37s/39$ and $27h-5s+8\eta\le216\tau$.
The oriented contraction-graph checker verifies the finite-sum
identities under dummy-index relabeling, in arbitrary dimension.
These formulas use actual entries of $f$; if one uses the coefficients
of $q$ instead, their common scale must be included before computing
$s,h,\eta,\tau$.

\subsection{The remaining block estimates}
The charges $+2,0,-2$ give the self-adjoint block pattern
\[
 S=\begin{pmatrix}D_+&L&0\\L^*&D_0&M\\0&M^*&D_-\end{pmatrix},\qquad
 \|D_+\|^2=\|D_-\|^2=\tau,\quad\|D_0\|^2=4\tau,
 \quad\|L\|^2=\|M\|^2=r=(h-6\tau)/4.
\]
The factor four follows by using the normalized symmetric mixed basis
$(e_i\otimes\bar e_j+\bar e_j\otimes e_i)/\sqrt2$.
Bounding each closed block walk by the product of its Hilbert--Schmidt
norms gives, with $T_j=\operatorname{tr}S^j$,
\[
\begin{split}
 T_3&\le10\tau^{3/2}+18r\sqrt\tau
          =\sqrt\tau(9h/2-17\tau),\\
 T_4&\le18\tau^2+56\tau r+8r^2
          =h^2/2+8h\tau-48\tau^2,\\
 \|\Pi_0(S^2)\|^2&\le18\tau^2+20\tau r+6r^2
          =3h^2/8+h\tau/2+3\tau^2/2,\\
 \|\Pi_{\pm2}(S^2)\|^2&\le36\tau r=9h\tau-54\tau^2.
\end{split}
\]
No positivity of $S$ is assumed. The low-range scalar argument also
uses $\|S\|_{\rm op}\le\sqrt{10/11}\|S\|$.
A short proof is to project $\operatorname{Sym}_4(B\otimes B)$, for
a unit real symmetric matrix $B$, onto charges $0,\pm2$. If $u$ is
the squared norm of its anti-commuting part, expansion of the three
pairings gives squared norm at most
$2/3+4u/3-11u^2/6\le10/11$.
Pairing with $S$ and taking the Rayleigh supremum proves the claim;
the detailed charge projection is in
\code{charge\_four\_operator\_bound.md}.

Write $G=G_0+G_2$, commuting and anti-commuting with the complex
structure $J$. Both $G$ and $J^*GJ=G_0-G_2$ are positive semidefinite.
Thus the particularly simple identity
\[
 0\le\operatorname{tr}(GJ^*GJ)=\|G_0\|^2-\|G_2\|^2=s-2\eta
\]
proves $\eta\le s/2$. Eigenvalues of $G_0$ have even real
multiplicity. It follows, for $V_3=\operatorname{tr}G^3$, that
\[
 0\le V_3\le(s+2\eta)\sqrt{(s-\eta)/2}.
\]
The components of $\operatorname{vec}G$ have norms squared
$\eta/2,s-\eta,\eta/2$. Hence
\[
 Q_G:=\langle\operatorname{vec}G,S\operatorname{vec}G\rangle
 \le(2s-\eta)\sqrt\tau+
           \sqrt{2\eta(s-\eta)(h-6\tau)}.
\]

Finally put $U=SF$, $D=FG$, and $k_\theta=\theta^2-2\theta$.
Orthogonal symmetrization preserves charge, so
$\langle C_j,S_j\rangle=\|S_j\|^2$ in each of the two sectors.
Direct substitution gives the exact nonnegative norms
\[
\begin{split}
 A_0&:=\|C_0-\theta S_0\|^2
      =2s/3-2h+(18+6k_\theta)\tau,\\
 A_2&:=\|C_2-\theta S_2\|^2
      =s/3+(2+k_\theta)h-(18+6k_\theta)\tau.
\end{split}
\]
Here the subscript $2$ includes both signs. Since $C-\theta S$ has
no charge-four component, Cauchy--Schwarz in these two sectors gives
\[
 \|U\|^2\le\theta T_3+
 \sqrt{A_0(3h^2/8+h\tau/2+3\tau^2/2)}+
 \sqrt{A_2(9h\tau-54\tau^2)}.
\]
This is valid for every real $\theta$, including $-54/125$.
Also $\|U\|^2\le\sqrt{sT_4}$ and
$\langle D,U\rangle\le\sqrt{V_3\|U\|^2}$.
These are all the tensor inequalities consumed by the scalar certificate.

\subsection{Cumulants, rational replay, and conversion}
Connected Wick contraction gives
\[
 \kappa_4(I_3(f))=1296(s+3h/2),
\]
\[
 \frac{\kappa_6(I_3(f))}{46656}
 =\frac{135}{2}T_3-20V_3+\frac{135}{2}Q_G
       +\frac{165}{2}\langle D,U\rangle+45\|U\|^2.
\]
The exact checker covers all 640 connected labelled cubic multigraphs
on six vertices, representing $11{,}314{,}080$ Wick pairings. The
negative $-20V_3$ is retained. For $s,h>0$ normalize by
\[
 y=\sqrt{h/s},\quad w=\eta/s,\quad
 u=\sqrt{V_3/s^{3/2}},\quad z=T_3/h^{3/2},\quad
 v=\|U\|^2/(h\sqrt s),\quad c=Q_G/(s\sqrt h).
\]
The real cumulant ratio is bounded by
\[
 \frac{(135/2)zy^3+45vy^2+(135/2)cy-20u^2
                 +(165/2)uy\sqrt v}{(1+3y^2/2)^{3/2}}.
\]
The archived checker with rational parameters
\code{split-moment-parameters.json} proves that this is at most
$(4045313/100000)10/(3\sqrt6)$, under the just-proved constraints.
It handles $y\le39/50$ algebraically and checks 100 affine-majorant
leaves, one derivative rectangle, and 17 high-range Bernstein leaves.
This is a fresh replay of that 118-leaf certificate; the supplied PDF
reports a different 59-leaf replay of the Lean polynomial data.
The zero cases follow directly from $F=0$ or $T_4=\|U\|^2=0$.

For $e=1$, circularity yields
\[
 \kappa_4(\sqrt2\Re q)=\tfrac32(m-2),\qquad
 \kappa_6(\sqrt2\Re q)=\tfrac52(\E|q|^6-9m+12).
\]
Substitution proves Theorem~\ref{thm:sixth-tensor-bridge-audit}; scaling
handles every $e$. Therefore the eleven-row radius theorem below is
unconditional within $\mathcal W_{2,1}$ as a mathematical \st{P,C}
result, while still requiring a Lean port.

\paragraph{Why the current broader Lean hypothesis cannot be erased.}
For independent real normalized Hermite cubics $A,B$, the variable
$q=(A+iB)/\sqrt2$ has moments $1,47,13506$ at orders $2,4,6$.
Its normalized sixth excess is
\[
 \frac{13506-9\cdot47+12}{(47-2)^{3/2}}
       =\frac{97}{\sqrt5}>40.45313.
\]
It is not in $\mathcal W_{2,1}$. The desired port strengthens the
premises to the correct bidegree, constructs the tensors, and then
applies the existing scalar theorem; it does not delete \code{h6}
from a theorem about arbitrary real third-chaos pairs.

\subsection{An exact vector witness and an improved elementary rate}
\begin{proposition}[\st{P}: vector transfer input]
\label{prop:exact-vector-transfer-audit}
Let $Z\in\mathbb C^n$ be standard proper Gaussian, $R=\|Z\|$,
$u_j=Z_j/R$ and $T_n=P_{1,0}-\epsilon P_{2,1}$, $0\le\epsilon\le1$.
Writing $\mu_n=\E R$, one has
\[
 \sum_j\langle u_j,T_nu_j\rangle
  =\frac{\mu_n^2}{n}\left(1-\frac{\epsilon}{4(n+1)}\right)
  >1-\frac{1+\epsilon}{4n}.
\]
Thus $n\ge(1+\epsilon)/(4\delta)$ suffices for vector value above
$1-\delta$. This proves the vector input on paper; it gives no scalar
operator-norm cap.
\end{proposition}
\begin{proof}
Phase integration and the second and fourth moments of the uniform
complex sphere give
\[
 P_{1,0}u_j=(\mu_n/n)Z_j,\qquad
 P_{2,1}u_j=-\frac{\mu_n}{2n(n+1)}Z_j(R^2-n-1).
\]
For the latter identity, pairing with each Wick basis monomial gives
$-\mu_n(\delta_{ja}\delta_{bc}+\delta_{jb}\delta_{ac})/(2n(n+1))$;
pairing $Z_j(R^2-n-1)$ gives the delta sum. Since
$\sum_j\E|Z_j(R^2-n-1)|^2=n(n+1)$, orthogonality proves the equality.

For a sharper elementary rate than the direct Cauchy--Schwarz bound,
use $\mu_{n+1}=(n+1/2)\mu_n/n$ and
$n^2/(n+1/2)\le\mu_n^2\le n$. The lower bound is Cauchy--Schwarz
applied to $\E R^2$ and $\E R\,\E R^3$.
Then $b_n=\mu_n^2/(n-1/4)$ tends to one and is strictly decreasing,
because
\[
 \frac{b_{n+1}}{b_n}
 =\frac{(n+1/2)^2(n-1/4)}{n^2(n+3/4)}
 =1-\frac1{16n^2(n+3/4)}<1.
\]
Hence $\mu_n^2>n-1/4$. Insert this in the exact payoff, and use
$\epsilon/(4(n+1))\le\epsilon/(4n)$ to obtain the stated rate.
\end{proof}

\paragraph{Git provenance.}
Commit \code{d0869cf} introduced $40.45313$ together with explicit
tensor assumptions. Commit \code{0969346} removed only the redundant
\code{sixth\_moment\_le\_nX} conversion wrapper, whose statement
already assumed the real cumulant bound. No removed unconditional
tensor proof was found. The distinction is external mathematics versus
the still-unconstructed bidegree bridge in Lean.
